\subsection{Terminal-based Interface Concepts}

This section describes the interfaces as such, and defines two interfaces that
are relevant to our solution: \glspl{cli} and \glspl{tui}, along with underlying
concepts. It also mentions the programs that leverage these interfaces.

\paragraph{Interface}
In computing, an interface is a shared boundary across which two or more
separate components of a computer system exchange information. The exchange can
be between software, computer hardware, peripheral devices, humans, and
combinations of these~\cite{hookway2014}.

\paragraph{Hardware terminal}
is a an electronic or electromechanical hardware device that can be used for
entering data into, and transcribing data from, a computer or a computing
system. The main functionality of terminals is to provide an elementary
text-based interface for user to interact with the computer
system~\cite{terminfo}.

\paragraph{Terminal emulator}
(later referred to as \textbf{Terminal})
is a software application that replicates the functionality of a traditional
hardware terminal as a software program. Here are some examples of terminal
emulators:
\begin{itemize}
  \item iTerm2 (macOS)
  \item Windows Terminal (Windows)
  \item Alacritty (Cross-platform)
\end{itemize}

\paragraph{Terminal application}
is a software program that runs within terminal emulators and utilizes their
capabilities to provide text-based user interfaces --- either \glspl{cli} or
\glspl{tui}. We will discuss it further in the following sections.

The communication between terminal applications and the terminal emulator is
typically done through standard input and output streams, which use file
descriptors as the underlying mechanism for data exchange. This restricts the
interaction to text-based commands and responses. Terminal applications can vary
from simple \gls{cli} utilities to complex \gls{tui}
programs~\cite{man_man,linfo_terminal}.

\paragraph{\gls{cli}}
The definition of \gls{cli} extends previous definition to the specific case of
command line interfaces, which are text-based interfaces between the user and
the computer system, where user inputs commands in the form of text, to
manipulate the system or perform specific tasks, and gets responses in the same
form.

\paragraph{Shell}
is special program (with \gls{cli}, \gls{tui}, as well as
\gls{gui} interfaces) that provides relatively wide and direct access on system
on which it runs.\ \gls{cli} shells are \gls{cli} programs that allow users to
execute commands and other \gls{cli} utilities~\cite{man_man,linfo_terminal}.

\paragraph{\gls{cli} utility}
is program executed from within the \gls{cli} environments. Some popular
examples of shells and \gls{cli} utilities are mentioned in Table~\ref{t:utils}.

\begin{table}
  \caption{Examples of shells and \gls{cli} utilities}\label{t:utils}
    \begin{tabularx}{\textwidth}{lXX}
    \toprule
    Name & System & Description \\
    \midrule
    Bash, Zsh
    & Linux, macOS
    & Popular Unix-like shells, compliant with POSIX standard.\\
    Cisco IOS Shell
    & Cisco devices
    & Shell used in Cisco network devices. \\
    bluetoothctl
    & Linux, macOS
    & \gls{cli} utility for managing Bluetooth devices.\\
    git
    & Cross-platform
    & \gls{cli} utility for version control using Git system. \\
    nmcli (NetworkManager)
    & Linux
    & Application for controlling network settings. \\
    \bottomrule
  \end{tabularx}
\end{table}

\paragraph{\gls{tui}}
is a type of user interface that is a further extension of \glspl{cli}. They
provide a more visually structured and intuitive way for users to interact with
applications by using libraries such as ncurses, which use the \textbf{terminal
capabilities} within emulators to create rich text-based
interfaces~\cite{linfo_terminal}.

\paragraph{Terminal capabilities}
are a special functionality of the terminal, which gives the application more
control over it's output. Such capabilities usually include the ability to move
the cursor, change text color, and rewrite or clean characters on the screen.
This is achieved by sending specific control sequences to the terminal emulator,
which then are interpreted to perform the desired actions. The examples of such
sequences are ANSI escape codes~\cite{terminfo,iso6429}. Some examples of
applications using \gls{tui} are mentioned in Table~\ref{t:utils2}.

\begin{table}
  \caption{Examples of \gls{tui} utilities and programs}\label{t:utils2}
    \begin{tabularx}{\textwidth}{lXX}
    \toprule
    Name & System & Description \\
    \midrule
    vi/vim/nvim
    & Linux, macOS, Windows
    & \gls{tui} visual
    text editors. \\
    mc (Midnight commander)
    & Linux, macOS, Windows
    & \gls{tui} file explorer.\\
    htop
    & Linux
    & Task manager.\\
    nmtui (NetworkManager)
    & Linux
    & Application for controlling network settings. \\
    \bottomrule
  \end{tabularx}
\end{table}


\subsection{\gls{llm} Concepts}

This section provides a brief overview of the aspects of \glspl{llm} and
embedded \glspl{llm}, relevant for this thesis. It describes the architecture,
types, limitations of \gls{llm}, the main protocols for the interaction, as well
as the advanced concepts.

\subsubsection{\glspl{llm}}

\paragraph{\glspl{llm}} refer to language models, typically based on transformer
architecture, characterized by having a large number of parameters, which are
trained on massive text data, such as \gls{gpt}~\cite{radford2018improving},
\gls{bert}~\cite{devlin2018bert}, \gls{lstm}~\cite{hochreiter1997lstm} or
Mamba~\cite{gu2023mamba}. \glspl{llm} exhibit strong capacities to understand
natural language and solve complex tasks (via text generation).

\paragraph{Language models} are mathematical models designed to predict and generate
sequences of text. They are crucial in artificial intelligence for performing
various tasks like language translation, summarization, and content creation.
Over time, they have evolved significantly, powered by advances in neural
networks~\cite{naveed2023}.

\subsubsection{\glspl{llm} Architecture Overview}

The architecture of \glspl{llm} can be divided into several main components:
\begin{itemize}
  \item \textbf{Tokenizer}: This component is responsible for providing the
    representation of the input text in a format that can be processed by
    the model. It breaks the input text into smaller units, to be able to
    represent it as vectors and then process it. The set of possible tokens
    is called the \textbf{vocabulary}, and can include words, subwords, or
    even individual characters. It is important to note that the vocabulary
    is fixed after training, which means that any new or rare words or
    combinations of characters (for example, command flags or arguments) may
    not be processed accurately or at all.
  \item \textbf{Embedding Layer}\label{sec:emb-layer}: This layer
    converts the tokenized input
    into dense vector representations, which capture the semantic meaning
    of each token. There are two main types of embeddings: \textbf{static}
    and \textbf{positional}. Static embeddings, such as Word2Vec or GloVe,
    assign a fixed vector to each token, regardless of the context in which
    it appears. Positional embeddings, on the other hand, generate
    different vectors for the same token based on its context in the
    input text, allowing for a more nuanced understanding of
    language~\cite{mikolov2013distributed,pennington2014glove,zhao2023}.
  \item \textbf{Transformer Layers}: These layers are the core of the \gls{llm}
    architecture, responsible for processing the input and generating the
    output. They consist of multiple attention~\cite{attention} heads and
    feedforward networks that allow the model to capture complex
    relationships between tokens and generate coherent and contextually
    relevant responses.
  \item \textbf{Output Layer}: This layer generates the final output of the
    model, which can be in the form of text, probabilities, or other formats
    depending on the task. It typically consists of a linear layer followed by
    a softmax activation function to produce a probability distribution over the
    vocabulary for text generation tasks.
\end{itemize}

In the original article, previously mentioned transformers
consist of two main parts~--- \textbf{Encoder} and \textbf{Decoder}.
The usage of the architecture based on attention implies that the
model generates the
output text sequentially, one token at a time, while attending to all
previously generated tokens. This can be a limitation for long prompts, as the
model can only attend to a limited number of previous tokens (this limit is
called \textbf{context length} or \textbf{window size}). Furthermore, the
longer the prompt, the more computationally expensive it becomes to generate
the output, as the model has to process all previous tokens for each new
token~\cite{attention}.

\textbf{Encoder}: The encoder maps an input sequence of symbols (tokens) to a
sequence of continuous representations. The encoder is composed of a stack of a
number of identical layers. Each layer has two sub-layers. The first is a
multi-head self-attention mechanism, and the second is a simple, position-wise
fully connected feed-forward network~\cite{attention}.

\textbf{Decoder}: The decoder then generates an output sequence of symbols
(tokens) one element at a time. At each step the model is auto-regressive,
consuming the previously generated symbols as additional input when generating
the next. The decoder is also composed of a stack of a number of identical
layers. In addition to the two sub-layers in each encoder layer, the decoder
inserts a third sub-layer, which performs multi-head attention over the output
of the encoder stack~\cite{attention}.

For the purpose of this thesis, it is enough to know that nowadays \glspl{llm}
are mainly based on Decoder-only architecture, which means they only use the
Decoder part of the original transformer architecture. Additionally,
Decoder-only models seem to be typically easier to train and fine-tune, as they
require less complex training data and objectives compared to Encoder-Decoder
models. That means that on the same numbers of training tokens, Decoder-only
models can achieve better performance than Encoder-Decoder models
~\cite{zhao2023,wei2023chainofthoughtpromptingelicitsreasoning}.

\subsubsection{\gls{moe} Extension}

While the architectures above define how information flows through the layers,
the \gls{moe} is an important extension that changes how the feedforward layers
operate, with the goal to increase scalability and efficiency.

In a traditional dense transformer, which is referred to as \textbf{Dense}, the
whole decoder is a monolith, meaning that while processing, every neuron is
used. The \gls{moe} architecture splits the massive, singular feedforward
network into multiple smaller, quasi-specialized (as everything is learned)
\textsl{expert} subnetworks and introduces a \textsl{router} (or gating network)
to orchestrate them. The router acts as a traffic controller, dynamically
selecting only a tiny subset of experts~--- typically one or two~--- to process
it.

When it comes to comparing \gls{moe} and dense models in scope of smaller
devices with certain limitations, \gls{moe} architectures introduce the changes
in two main aspects compared to dense models: computational efficiency and
memory usage.
\begin{itemize}
  \item \textbf{Computational Efficiency:} \gls{moe} models are more
    computationally
    lighter as the selective activation leads to significant reductions in
    computation time and energy consumption. However, such step makes the
    model less intelligent, especially on smaller scales. In contrary, a
    dense model involves every parameter in the computation for each token,
    leading to higher computational costs, but also more consistent performance.
  \item \textbf{Memory Usage:} The major drawback of \gls{moe} models
    is their high
    memory requirements. Since all experts must be loaded into memory to allow
    the router to select them dynamically, the total memory footprint is
    much larger than that of a dense model. Additionally, such architecture
    allows for knowledge duplication, thus less efficient memory usage.
\end{itemize}
Overall, \gls{moe} models tend to be more efficient in terms of computation, but
have worse performance over dense models with same number of parameters.
However, if VRAM is not a constraint, MoE models can achieve better performance
than dense models within the same computational budget.

\subsubsection{Embedding \gls{llm}}\label{sec:embedding-models}

\paragraph{Embedding model}
is a type of neural network architecture that aims to extract features from
complex data (such as text, audio, video, or other types of data), mapping them
to numerical representations (embeddings). These embeddings are much better for
the machine to process it in comparison to the raw text, as they capture the
semantic meaning of the text and can be used in various
applications~\cite{nie2025textembeddingmeetslarge}.

\paragraph{Embedding \gls{llm}} 
is a embedding model, based on the \textbf{transformer decoder architecture}
that is specifically designed to generate dense vector representations of
content. These embeddings capture the semantic meaning of the content, allowing
for tasks such as semantic search, clustering, and other applications that
require understanding of the content's meaning.

The advantage of using embedding \gls{llm} in comparison to traditional
embedding models is that they can capture more complex relationships between
tokens and provide a more accurate representation of the content's meaning. This
is because embedding \gls{llm} are typically based on transformer architectures,
which use attention mechanisms to weigh the importance of different tokens in
the input text, allowing them to capture context and relationships between
tokens more effectively than traditional embedding
models~\cite{zhang2024systematicsurveytextsummarization,khan2026retrieval}.

\subsection{Advanced \gls{llm} Concepts}

This section describes two most important advanced concepts related to
\glspl{llm} that are relevant for this thesis: \textbf{\gls{llm} Agents} and
\textbf{\gls{rag}}.

\subsubsection{\gls{llm} AI Agent Systems}\label{agent}

The evolution of \glspl{llm} has progressed beyond simple text, and
\glspl{llm} have been increasingly integrated into more complex systems and
applications, often in context of \textbf{AI Agents}, which utilize \glspl{llm}
with the capacity to plan, act, and interact with external
environments~\cite{naveed2023,zhao2023}.

\paragraph{AI agents} are autonomous entities, capable of planning,
decision-making, and performing actions to achieve complex
goals~\cite{naveed2023,zhao2023}.

We can define the workflow of \gls{llm} Agents ~\cite{nvidia-llm-agent} as
pseudocode in Listing~\ref{alg:agent}.

\begin{algorithm}[t]
\caption{Pseudocode for \gls{llm} Agent workflow}
\label{alg:agent}
\begin{algorithmic}[1]
\STATE \textbf{function} \textsc{LLM\_Agent}(\textit{task}, \textit{context})
    \STATE \quad $\text{results} \gets \text{empty list}$ \COMMENT{Initialize results list}
    \STATE \quad $\textit{subtasks} \gets \textit{LLM}.\text{decompose\_task}(\textit{task})$ \COMMENT{Break down main task into sub-tasks}

    \FOR{\textbf{each} $\textit{subtask}$ \textbf{in} $\textit{subtasks}$}
        \IF{$\textit{subtask}.\textit{type} = \text{tool}$}
            \STATE $\textit{result} \gets \text{use\_tool}(\textit{subtask}.\textit{tool}, \textit{subtask}.\textit{tool\_input})$
        \ELSIF{$\textit{subtask}.\textit{type} = \text{further\_decomposition}$}
            \STATE \COMMENT{Recursive call for further decomposition}
            \STATE $\textit{result} \gets \text{AI\_Agent}(\textit{subtask}.\textit{task}, \textit{context})$
        \ELSIF{$\textit{subtask}.\textit{type} = \text{generate}$}
            \STATE $\textit{result} \gets \textit{LLM}.\text{generate\_response}(\textit{subtask}.\textit{task})$
        \ENDIF
        \STATE $\text{results}.\text{append}(\textit{result})$
    \ENDFOR

    \STATE \quad \textbf{return} \textit{LLM}.\text{summarize\_results}(\textit{results})
\STATE \textbf{end function}
\end{algorithmic}
\end{algorithm}

AI Agent architecture brings major improvements, compared to traditional
\glspl{llm}:
\begin{itemize}
  \item \textbf{Hierarchical Planning:} AI Agents can break down complex tasks into
    smaller, more manageable sub-tasks, allowing for more efficient problem-solving.
  \item \textbf{Tool Use and Abstraction:} AI Agents can leverage external tools,
    RAG pipelines or APIs to retrieve information or perform actions
    beyond text generation, while keeping only the relevant information.
  \item \textbf{Iterative Generation:} AI Agents can generate multiple steps of
    reasoning or actions iteratively, allowing them to refine their outputs
    and adapt to new information.
\end{itemize}
In essence, AI Agents perform tasks, dividing complex objectives into
manageable steps, and iteratively prompting themselves to execute
these steps using various tools.


\subsubsection{Agent tools and MCP Protocol}

\gls{mcp} protocol is a protocol, that standartizes the interaction between
agent and tools, allowing simpler software integration and development for AI
agents~\cite{anthropic2024mcp}. While not used in this thesis, it is a relevant
and important concept in the field of \glspl{llm}.


\subsubsection{Retrieval-Augmented Generation (RAG)}

\textbf{\gls{rag}} is a technique that allows \glspl{llm} to utilize external
knowledge bases, significantly reducing hallucinations and increasing
accuracy~\cite{zhao2023}. It reuses the semantic meaning extraction capabilities
of \textbf{embedding \gls{llm}} and the text generation and summarization
capabilities of \textbf{\gls{llm}} to retrieve information from external
sources. We define the workflow of \gls{rag} in Listing~\ref{alg:rag}.

\begin{algorithm}[t]
\caption{Pseudocode for \gls{rag} workflow}
\label{alg:rag}
\begin{algorithmic}[1]
\STATE \textbf{function} \textsc{LLM\_RAG\_generate}(\textit{user\_query}, \textit{LLM}, \textit{document\_database})
    \STATE \quad \COMMENT{Step 1: Retrieve relevant documents from the database}
    \STATE \quad $\textit{relevant\_documents} \gets \text{retrieve\_documents}(\textit{document\_database}, \textit{user\_query})$

    \STATE \quad \COMMENT{Step 2: Inject retrieved documents into the input prompt}
    \STATE \quad $\textit{augmented\_prompt} \gets \text{construct\_prompt}(\textit{user\_query}, \textit{relevant\_documents})$

    \STATE \quad \COMMENT{Step 3: Pass the augmented prompt to the LLM}
    \STATE \quad $\textit{llm\_response} \gets \textit{LLM}.\text{generate\_response}(\textit{augmented\_prompt})$

    \STATE \quad \COMMENT{Step 4: Return the final response}
    \STATE \quad \textbf{return} $\textit{llm\_response}$
\STATE \textbf{end function}
\end{algorithmic}
\end{algorithm}

This technique involves document retrieval, the best general approach for which
is to use \textbf{semantic document search and retrieval}
(see~\ref{emb:semvec}). By this, we mean the set of techniques of indexing and
searching documents based on their semantic representations. Therefore, there
are two main components for the retrieval system: document representation
generator (embedding model, mentioned in Section~\ref{sec:embedding-models}) and
storage with searching capabilities suited for the task
(see~\ref{sec:vector-db}).

\paragraph{Different Granularities In Retrieval}

\gls{rag} systems can retrieve information at multiple granularities, for
example:

\begin{itemize}
  \item \textbf{Document-level retrieval}
  \item \textbf{Passage-level retrieval}
  \item \textbf{Sentence-level retrieval}
\end{itemize}

The choice of granularity defines the smallest unit, for which the (vector)
representation is generated. This affects the computational load for the
\gls{rag} system for the creation and retrieval, or how much different units can
be fit into \gls{llm} context window or answer size.

%  Level-1
%
% Explicit Facts: These queries are asking about explicit facts
% directly present in the given data without requiring any additional
% reasoning. This is the simplest form of query, where the model’s
% task is primarily to locate and extract the relevant information.
% For example, "Where will the 2024 Summer Olympics be held?" targets
% a fact contained in the external data.
% Level-2
%
% Implicit Facts: These queries ask about implicit facts in the data,
% which are not immediately obvious and may require some level of
% common sense reasoning or basic logical deductions. The necessary
% information might be spread across multiple segments or require
% simple inferencing. For instance, the question "What is the
% majority party now in the country where Canberra is located?" can
% be answered by combining the fact that Canberra is in Australia
% with the information about the current majority party in Australia.
% Level-3
%
% Interpretable Rationales: These queries demand not only a grasp of
% the factual content but also the capacity to comprehend and apply
% domain-specific rationales that are integral to the data’s context.
% These rationales are often explicitly provided in external
% resources and is typically not present or rarely encountered during
% the pre-training phase of a general large language model. For
% example, in the realm of pharmaceuticals, an LLM must interpret FDA
% Guidance1 documents—which represent the FDA’s current thinking—to
% evaluate whether a specific drug application adheres to regulatory
% requirements. Similarly, in customer support scenarios, the LLM
% must navigate the intricacies of a predefined workflow to process
% user inquiries effectively. In the medical field, many diagnostic
% manuals provide authoritative and standardized diagnostic criteria,
% such as management guidelines for patients with acute chest pain
% [14]. By effectively following these external rationales, it is
% possible to develop a specialized LLM expert system for managing
% chest pain. This involves understanding the procedural steps and
% decision trees that guide a support agent’s interactions with
% customers, ensuring responses are not only accurate but also comply
% with the company’s service standards and protocols.
% Level-4
%
% Hidden Rationales: This category of queries delves into the more
% challenging realm where the rationales are not explicitly
% documented but must be inferred from patterns and outcomes observed
% in external data. The hidden rationales here refer not only to the
% implicit reasoning chains and logical relationships, but also to
% the inherently challenging and non-trivial task of identifying and
% extracting the external rationales required for each specific
% query. In IT operational scenarios, for example, a cloud operations
% team may have addressed numerous incidents in the past, each with
% its own unique set of circumstances and resolutions. The LLM must
% be adept at mining this rich repository of tacit knowledge to
% discern the implicit strategies and decision-making processes that
% were successful. Similarly, in software development, the debugging
% history of previous bugs can provide a wealth of implicit insights.
% While the step-by-step rationale for each debugging decision may
% not be systematically recorded, the LLM must be capable of
% extracting the underlying principles that guided those decisions.
% By synthesizing these hidden rationales, the LLM can generate
% responses that are not only accurate but also reflective of the
% unspoken expertise and problem-solving approaches that have been
% honed over time by experienced professionals.

\paragraph{Query Types}
We can generally classify queries that user can ask to the \gls{rag} system into
four main types, defined in the original
article~\cite{zhao2024retrievalaugmentedgenerationrag}:

\begin{itemize}
  \item \textbf{Level-1: Explicit Facts:} These queries are asking
    about explicit facts
    directly present in the given data without requiring any additional
    reasoning. This is the simplest form of query, where the model’s task is
    primarily to locate and extract the relevant information.
  \item \textbf{Level-2: Implicit Facts:} These queries ask about implicit
    facts in the data, usually searching for the information in multiple
    segments and require some level of common sense reasoning or basic logical
    deductions.
  \item \textbf{Level-3: Interpretable Rationales:} To answer these queries,
    the model needs to not only retrieve and understand the facts in the
    database, but also suit the answer to the specific domain, by
    understanding and applying the domain-specific rationales.
  \item \textbf{Level-4: Hidden Rationales:} The answer to these queries is
    usually not explicitly documented, but must be deduced from patterns and
    outcomes found in external data.
\end{itemize}

With such definition of query types, it is possible to map the responsibilities
of the parts of the \gls{rag} system to the query types.

Level-1 queries can be mostly handled by the retrieval component, with the role
of the \gls{llm} being to filter the relevant information from the retrieved
documents.\\
Level-2 queries require the \gls{llm} to perform some reasoning and inference,
to deduct Level-1 queries from the Level-2 to ask the retrieval component, and
then combine the retrieved information to produce the final answer.\\
Level-3 queries require the \gls{llm} to understand the domain-specific
rationales, usually to be pre-trained on the documents containing these
rationales, and then apply them to the retrieved information to produce the
final answer. Along with Level-4 it requires advanced reasoning capabilities.\\

In this thesis, we will be focusing on the Level-1 and Level-2 queries, as they
are more relevant for the task of choosing the right document based on the user
query, and the application of the system in the context of operating on
manuals usually requires the retrieval of explicit and implicit facts, rather
than the deduction of hidden rationales.

\paragraph{Common RAG Testing Techniques}\label{sec:rag-eval}

This part describes most common RAG Testing Techniques, essential for the
evaluation of the performance of \gls{rag} systems~\cite{rag-evaluation,
langsmith2026rag}.

\paragraph{Precision@k}

\textbf{Precision@k} measures the proportion of relevant documents or passages
that appear within the top-k retrieved results. It is defined as:

\[
  \text{Precision@k} = \frac{\text{number of relevant documents in
  top-k}}{\text{k}}
\]

This metric directly assesses the quality of the retrieval component. For
example, Precision@1 measures whether the most relevant document is ranked
first, while Precision@5 measures whether relevant documents appear within the
top 5 results. Higher precision at low k values (e.g., k=1) is generally
preferred, as it indicates the system's top results are highly relevant.

\paragraph{Recall@k}

\textbf{Recall@k} measures the proportion of all relevant documents that are
successfully retrieved within the top-k results:

\[
  \text{Recall@k} = \frac{\text{number of relevant documents
  retrieved in top-k}}{\text{total number of relevant documents}}
\]

Recall provides a complementary perspective to precision: a system can achieve
high precision by returning few results, but recall indicates whether the system
finds a significant portion of all relevant documents. Recall@k is particularly
important when comprehensive coverage is needed.

\paragraph{Mean Reciprocal Rank (MRR)}

\textbf{Mean Reciprocal Rank} measures the average rank of the first relevant
result across all queries:

\[
  \text{MRR} = \frac{1}{|Q|} \sum_{i=1}^{|Q|} \frac{1}{\text{rank}_i}
\]

where \textsl{rank\textsubscript{i}} is the position of the first relevant
document for query $i$, and |Q| is the total number of queries. MRR rewards
systems that find relevant documents early in the ranking and is particularly
useful when users typically only examine the top few results.

% \paragraph{Normalized Discounted Cumulative Gain (NDCG)}
%
% \textbf{\gls{ndcg}} combines relevance judgments with
% ranking position to account for the fact that lower-ranked relevant documents
% are less valuable than higher-ranked ones:
%
% \[
%   \text{DCG} = \sum_{i=1}^{k} \frac{\text{rel}_i}{\log_2(i+1)}
% \]
%
% where \textsl{rel\textsubscript{i}} is the relevance score of the document at
% position $i$.\ \gls{ndcg} normalizes this by the ideal ranking (IDCG) to produce a
% score between 0 and 1. NDCG is valuable for assessing ranking quality when there
% are multiple levels of relevance (e.g., highly relevant, partially relevant, not
% relevant).
%
% \paragraph{F1 Score for Retrieval}
%
% For binary relevance judgments (relevant or not relevant), the \textbf{F1 score}
% combines precision and recall:
%
% \[
%   \text{F1} = 2 \times \frac{\text{Precision} \times
%   \text{Recall}}{\text{Precision} + \text{Recall}}
% \]
%
% The F1 score is useful for balancing the trade-off between precision and recall
% when a single metric is needed to optimize system performance.

\subsubsection{LLM Fine-tuning and Training}

\paragraph{Training} is the process of teaching a \gls{llm} to understand and
generate text with large datasets, which involves creating a model from scratch.

\paragraph{Fine-tuning} is the process of taking a pre-trained \gls{llm} and
adapting it to a specific task by training it on a smaller, task-specific
dataset. This allows the model to learn task-specific patterns and improve its
performance on the target task while keeping low training effort and
costs~\cite{zhao2023}.

While not used in this thesis, it is important to mention these techniques as
relevant concepts.

\subsection{Document Representation And Retrieval}

This section provides the background on the techniques for representing
documents in a way that allows for accurate retrieval. It also describes the
techniques for generating the semantic vector representations of documents, as
well as the storage and retrieval of these representations using vector
databases.

\subsubsection{Document Representation}\label{emb:semvec}

\paragraph{Semantic representation of documents} is the process of converting
documents into a format that captures their meaning. In scope of this thesis, we
are interested in the techniques that create numerical representations of
documents, which are much better for the machine to process. Currently, the
semantic representation of documents is generally categorized into two primary
groups~\cite{khan2026retrieval}:

% TODO: citations
\begin{itemize}
  \item \textbf{Sparse vector representations (sparse retrieval):} This method
      mostly relies on traditional keyword-based approaches, such as \gls{bow},
        one-hot encoding, \gls{lsindexing} or \gls{tfidf} to retrieve documents
        based on the presence or quantity of specific terms. It is
        computationally efficient and works well for simple queries, but it may
        struggle with complex queries or those that require understanding of
        context and semantics~\cite{khan2026retrieval}.

        \textsl{The common example of sparse retrieval is \gls{bow} model, which
        represents documents as vectors of word counts. It needs the dictionary
        of all possible words (or subset of them) to be defined. Then, the size
        of the vector is equal to the size of the dictionary, and each element
        position corresponds to the word in the dictionary.}

  \item \textbf{Dense vector representations (dense retrieval):} This method
      uses dense vector representations of documents and queries to calculate
        similarity. It typically involves training a neural network to generate
        embeddings that capture the semantic meaning of the text. The similarity
        between the query and documents is then calculated using cosine
        similarity or other distance metrics. Dense retrieval can be more
        effective for complex queries and can capture nuances in language, but
        it may require more computational resources compared to traditional
        sparse retrieval~\cite{khan2026retrieval}.
\end{itemize}

The original article also mentions \textbf{hybrid retrieval techniques}, that
aim to combine the computational efficiency of sparse methods with the ability
of dense representations to alleviate the vocabulary mismatch problem described
in the article. These techniques typically utilize \textbf{document expansion}
to bridge the gap between exact lexical matches and semantic meaning.
\textsl{For example, the SPLADE model expands the sparse representation of a
document or query with semantically related terms, allowing for better retrieval
performance even when exact keywords are missing.} The thesis is not going to
focus on the hybrid retrieval techniques, but it is important to acknowledge the
existence of such techniques, as they can be used in the future research to
further improve the retrieval performance of the system~\cite{khan2026retrieval}.

For the purpose of this research, we will focus on dense retrieval, as it is
more suitable for complex queries and can capture the semantic meaning of the
documents, which is more relevant for the task of choosing the right document
based on the user query.

% TODO: why going to dense (natlang, better acc)
\paragraph{Dense representation vector generation}
is an important aspect of the representation of documents. As the embedding
layer~\ref{sec:emb-layer} of either dense or sparse retrieval represents the
document as a collection of vectors, one for each token, it is important to have
the techniques to generate the single vector that represents this whole
collection.

\begin{itemize}
  \item \textbf{Classical Approaches:} These techniques include
    methods like averaging, max pooling and other simple
    aggregation methods to combine the token embeddings into a
    single vector. These methods are computationally efficient
    and easy to implement, but they may not capture the full
    semantic meaning of the document, but rather features like
    frequency of words or presence and relationships based on it,
    but not the context.
  \item \textbf{Neural Network-based Approaches:} These techniques involve
    training a neural network, such as a transformer (embedding
    model, mentioned in~\ref{sec:embedding-models}) or recurrent
    neural network, to generate a single vector representation for
    the entire document. The possible approaches are:
    \begin{itemize}
      \item \textbf{Multi-head attention-based summarization:} This method
        uses a transformer architecture retrieve weights for each token
        based on its position, relationships with other tokens and
        context, normalizes them, and then combines the token
        embeddings into a single vector using these weights. This
        approach can capture complex relationships between tokens and
        provide a more accurate representation of the document's
        meaning~\cite{sharma2024unmacapsumtunifiedmultiheadattentiondriven}.
      \item \textbf{[CLS] token-based summarization:} This method uses a
        special token (e.g., [CLS] in \gls{bert}) that is added to the
        beginning of the document. The transformer is trained to update
        the values of this token based on the context of the entire
        document, and the final value of this token is used as the
        summary vector for the
        document~\cite{behrendt2025maxpoolbertenhancingbertclassification,devlin2018bert}.
    \end{itemize}
    \textsl{It is important to note that to increase the accuracy of the
    summarization, it is often necessary to have single-topic documents.}
\end{itemize}

\subsubsection{Document Storage and Retrieval}

When document representations are generated, they need to be stored in a way
that allows for efficient retrieval based on similarity. For chosen purpose, we
focus on the storage and retrieval of dense vector representations. For this
purpose, \textbf{vector databases} are used.

\paragraph{Vector databases}\label{sec:vector-db} are specialized databases
designed to store and manage high-dimensional vector data and perform efficient
similarity search based on different metrics (e.g., cosine similarity, L2 or
Euclidean distance). Examples of vector databases include
\gls{faiss}~\cite{faiss2026}, \gls{annoy}~\cite{annoy2026}, and
Milvus~\cite{milvus2026}. Simplest way to use vector database is to store the
vector representations of documents along with corresponding metadata (e.g.,
document identifier, title), and then retrieve the most similar documents based
on the similarity between the vector, generated from the user query, and the
stored document vectors.

% TODO: more details on vector databases, how they work, how to use them, etc.


\subsection{Linux Manual Pages}

In this thesis and implementation, we focus on processing \textbf{Linux man
manual pages} (shortened \textbf{man pages} or \textbf{manual pages}) as data
for testing and validation and by the term \textbf{document} we will mean a man
page. In this part we provide the description of the main conventions for man
pages, which are relevant for the solution.

\texttt{man} (software application) is the system's manual pager. Each page
argument given to man is normally the name of a program, utility or function.
The manual page associated with each of these arguments is then found and
displayed. A section, if provided, will direct man to look only in that section
of the manual~\cite{man_man}.

Conventionally, manual pages are divided into several groups, and each manual
page consists of several sections. Conventional groups and section names are
listed in Appendix~\ref{a:man}.
The most important sections for the purpose of this thesis are
\textbf{DESCRIPTION}, which provides a detailed description of the program, its
functionality and usage, and \textbf{OPTIONS}, which lists the flags and
arguments that can be used with the program, along with their descriptions.
